% !TEX program = xelatex
% !TEX options = -synctex=1 -interaction=nonstopmode -file-line-error
% example-clinical.tex
% Demonstrates rossbeamerclinical features: inclusion/exclusion,
% endpoint cards, protocol timeline, CONSORT flow, safety card.
\documentclass[palette=clinical]{rossbeamerclinical}

\title{CLARITY-AD: Phase III Results}
\subtitle{Lecanemab in Early Alzheimer's Disease}
\author{Dr.~Elena Vasquez}
\institute{Memory Disorders Research Centre}
\date{November 2026}
\trialid{NCT01234567}
\sponsor{Biogen / Eisai}

% Dummy bib
\begin{filecontents}[overwrite]{\jobname.bib}
@article{clarity2023,
  author  = {van Dyck, C. H. and Swanson, C. J.},
  title   = {Lecanemab in Early {Alzheimer's} Disease},
  journal = {NEJM},
  year    = {2023},
  volume  = {388},
  pages   = {9--21},
}
\end{filecontents}
\addbibresource{\jobname.bib}

\begin{document}

% --- Title ---
\begin{frame}[plain]
  \titlepage
\end{frame}

% --- Inclusion / Exclusion ---
\begin{frame}{Eligibility Criteria}
  \inclusionexclusion{%
    \item Age 50--90 years
    \item Clinical diagnosis of MCI or mild AD dementia
    \item Amyloid PET positive or CSF A$\beta$42 $<$\,600\,pg/mL
    \item MMSE 22--30
    \item Stable concomitant medications $\geq$\,4 weeks
  }{%
    \item Known non-AD dementia
    \item Clinically significant cerebrovascular disease
    \item Contraindication to MRI
    \item Active malignancy within 5 years
    \item Prior anti-amyloid immunotherapy
  }
\end{frame}

% --- Protocol Timeline ---
\begin{frame}{Study Schedule}
  \begin{protocoltimeline}
    \visitpoint{0.5}{Screening}{Consent,\\MRI, PET}
    \visitpoint{3}{Baseline}{Cognitive\\battery}
    \visitpoint{5.5}{Week 12}{Safety\\MRI}
    \visitpoint{8}{Week 26}{Interim\\analysis}
    \visitpoint{10.5}{Week 52}{Primary\\endpoint}
  \end{protocoltimeline}
\end{frame}

% --- CONSORT Flow ---
\begin{frame}{CONSORT Flow}
  \centering
  \flowbox{Assessed for eligibility\\($n = 2{,}840$)}
  \flowarrow
  \begin{columns}[c]
    \begin{column}{0.55\textwidth}\centering
      \flowbox{Randomised\\($n = 1{,}795$)}
    \end{column}
    \begin{column}{0.40\textwidth}\centering
      \flowboxexcluded{Excluded ($n = 1{,}045$)\\Screen failure: 812\\Withdrew consent: 233}
    \end{column}
  \end{columns}
  \flowarrow
  \randomise{1{,}795}{Lecanemab ($n = 898$), Placebo ($n = 897$)}
\end{frame}

% --- Endpoint Cards ---
\begin{frame}{Primary and Key Secondary Endpoints}
  \endpointcard{Primary: CDR-SB Change at 18 Months}
    {$-0.45$}{$-0.67,\;-0.23$}{$< .001$}
  \vskip6pt
  \endpointcard{Secondary: ADAS-Cog14 Change}
    {$-1.44$}{$-2.92,\;0.04$}{.056}
  \vskip6pt
  \endpointcard{Secondary: ADCS-MCI-ADL Change}
    {$-2.0$}{$-3.5,\;-0.5$}{.009}
\end{frame}

% --- Safety Card ---
\begin{frame}{Safety Summary}
  \safetycard{ARIA-E (any grade)}{12.6\%}{1.7\%}{$<.001$}
  \vskip4pt
  \safetycard{ARIA-H microhaemorrhage}{14.0\%}{7.7\%}{$<.001$}
  \vskip4pt
  \safetycard{Infusion-related reactions}{26.4\%}{7.4\%}{$<.001$}
  \vskip4pt
  \safetycard{Serious adverse events}{14.0\%}{11.3\%}{.10}
  \vskip6pt
  \NNT{6} for CDR-SB responder\qquad \NNH{9} for symptomatic ARIA-E
\end{frame}

% --- Finding ---
\begin{frame}{Interpretation}
  \begin{findingbox}
    Lecanemab reduced clinical decline on CDR-SB by 27\% compared to
    placebo at 18 months in participants with early AD confirmed by
    amyloid biomarkers.
  \end{findingbox}
  \vskip6pt
  \begin{limitationbox}
    ARIA remains a significant safety signal. ApoE4 homozygotes had
    higher rates of ARIA-E (35\% vs 10\% for non-carriers).
  \end{limitationbox}
\end{frame}

\conclusionslide{Thank You}{e.vasquez@memorycentre.org}

\end{document}
